\section{管理建议与推广分析}

\subsection{日常排班建议}

实际运行时可在每天开始前读取最新预测，重新求解当天的5个班次。由于模型的输入只有逐小时到货量，预测更新后不需要修改约束结构。对第12天这类高峰日，应提前确认581名问题二所需人员；对低峰日，可根据生成的班次人数减少到岗规模，同时保持16时服务承诺。

\subsection{机动人员与安全余量}

问题三的581名固定员工每天并非全部到岗，具体到岗人数由出勤矩阵决定；第12天全员到岗以覆盖峰值。效率下降5\%后峰值增至611人，超过基准峰值30人，说明固定人员方案对效率波动敏感。若实际业务存在迟到或缺勤，建议在581人的基础上增加临时工或设置跨班次机动池，并在模型中将最低覆盖改为$r_d+s_d$。

\subsection{多目标扩展}

当前目标只最小化人数，未区分早班、晚班和夜班成本。推广时可使用
\begin{equation}
\min\sum_{d,s}c_s y_{ds}+\lambda\sum_{d,h}b_{dh}+\mu\sum_{i,d}|w_{i,d+1}-w_{id}|,
\end{equation}
分别控制用工成本、积压惩罚和出勤切换。若加入公平性约束，可以限制不同工人的夜班数量差异；若到货量为随机变量，可用多个情景共用班次决策，或采用鲁棒容量约束。

\subsection{模型推广边界}

模型适用于处理能力可加、班次长度固定且货物可排队的分拣系统。对于必须按订单优先级处理、不能混合排队或存在多类技能工人的中心，需要把货物类别、工人技能和设备资源加入流量网络。此时本文的库存守恒思想仍然适用，但单一产能系数应替换为多资源能力矩阵。

\section{全文总结}

本文从原始附件出发建立了可执行的动态排班系统，并将理论下界、整数求解、逐小时库存、月度人员约束和敏感性分析串联起来。最终结果不是孤立的一个人数，而是包括班次、人数、处理量、库存、低效率工时、工人日期出勤和班次分配在内的完整方案。该结果既满足训练题目的三项要求，也为将模型迁移到真实分拣中心提供了接口。
